\subsubsection*{引理 5.3.}
\underline{设 $X$ 为维数 $n$ 的光滑连通拟射影簇, $M$ 为 $X$ 上有限秩 $a$ 的局部自由层. 以 $\xi$ 表示 $\mathbb P(M)$ 上的可逆层 $\mathcal O_{\mathbb P(M)}(1)$, 设 $\pi:\mathbb P(M)\to X$ 为结构态射. 回顾 $\operatorname{CH}(\mathbb P(M))$ 是(经 $\pi^*$ 作用的)自由 $\operatorname{CH}(X)$-模, 基为 $1,\xi,\ldots,\xi^{a-1}$, 其中 $\xi$ 满足}
\begin{equation}\tag{5.3.1}
\xi^a-C_1\xi^{a-1}+C_2\xi^{a-2}-C_3\xi^{a-3}+\cdots=0,
\end{equation}
\underline{$C_i$ 为 $M$ 在 $\operatorname{CH}(X)$ 中的 Chern 类. 置 $C(M)=1+C_1+C_2+\cdots$, 即 $M$ 的全 Chern 类; 它是 $\operatorname{CH}(X)$ 的可逆元. 则:}

\begin{enumerate}
\item[(i)] \underline{有}
\begin{equation}\tag{5.3.2}
\pi_*\!\left(\frac{\xi^{a-1}}{1+\xi}\right)
=\frac{1}{C(M)};
\end{equation}
\item[(ii)] \underline{若此外 $X$ 固有, 设 $F(T)\in\operatorname{CH}(X)[T]$ 为多项式, 其中 $T^i$ 的系数为 $n+a-1-i$ 次齐次元素. 因而 $F(\xi)$ 在 $\operatorname{CH}(\mathbb P(M))$ 中具有最高次数 $n+a-1=\dim\mathbb P(M)$. 则}
\begin{equation}\tag{5.3.3}
\int_{\mathbb P(M)}F(\xi)
=(-1)^{a-1}\int_X\frac{F(-1)}{C(M)}.
\end{equation}
\end{enumerate}
\EditorialFootnote{积分的适用范围: 5.2 将积分定义为零维 Chow 类的次数. 对一般非固有拟射影簇, 这个次数不在有理等价下良定; 因此本引理的 (ii) 需加固有性, 而 (i) 的射影丛推出恒等式不需此附加假设. 法文原书与英文底本未在 (ii) 中补述这一条件. 本讲实际应用于射影簇 $X$, 所以该应用满足条件.}

\noindent\underline{证明.}
首先由 (5.3.2) 推出 (5.3.3). 写成
\begin{equation}\tag{5.3.4}
\frac{1}{C(M)}=\sum_{j\geq0}b_j,\qquad
b_0=1,\qquad
b_j\in\operatorname{CH}^j(X),\qquad
b_j=0\ \text{对 }j<0.
\end{equation}
则 (5.3.2) 等价于
\begin{equation}\tag{5.3.5}
\pi_*(\xi^{a-1+j})=(-1)^jb_j
\qquad\text{当 }a-1+j\geq0,
\end{equation}
因为由维数可知, 当 $q<a-1$ 时有 $\pi_*(\xi^q)=0$.

由线性性, 只需在 $F(T)=uT^{a-1+j}$ 时证明 (5.3.3), 其中 $u\in\operatorname{CH}^{n-j}(X)$ 为齐次元素. 由 (5.3.5) 及投影公式,
\begin{equation}\tag{5.3.6}
\begin{aligned}
\int_{\mathbb P(M)}F(\xi)
&=\int_{\mathbb P(M)}\pi^*(u)\xi^{a-1+j}\\
&=\int_Xu\,\pi_*(\xi^{a-1+j})
=(-1)^j\int_Xu\,b_j.
\end{aligned}
\end{equation}
另一方面,
\begin{equation}\tag{5.3.7}
\begin{aligned}
(-1)^{a-1}\int_X\frac{F(-1)}{C(M)}
&=(-1)^{a-1}\int_X
  \frac{(-1)^{a-1+j}u}{C(M)}\\
&=(-1)^j\int_X\sum_i u\,b_i.
\end{aligned}
\end{equation}
由于 $\deg u=n-j$ 且 $\deg b_i=i$, 有
\[
\int_Xu\,b_i=0\qquad\text{对 }i\ne j.
\]
故 (5.3.7) 化为
\begin{equation}\tag{5.3.8}
(-1)^{a-1}\int_X\frac{F(-1)}{C(M)}
=(-1)^j\int_Xu\,b_j,
\end{equation}
这就证明了所断言的蕴含.

尚需证明 (5.3.2), 或等价地证明 (5.3.5).

